\documentclass[10pt,a4paper]{amsart}
\usepackage[margin=19mm]{geometry}
\usepackage{amsmath,amssymb,mathtools}
\usepackage[scr=rsfs,cal=euler]{mathalfa}
\usepackage{xfrac}
\usepackage[unicode,hidelinks,bookmarks=false]{hyperref}
\newcommand{\ie}{i.e.\ }
\newcommand{\cf}{cf.\ }
\newcommand{\resp}{respectively\ }
\newcommand{\Subsection}{\subsection}
\providecommand{\nobreakdash}{\nobreak\hbox{-}}
\setlength{\emergencystretch}{2em}
\multlinegap0pt
\usepackage{amssymb}
\usepackage{mathtools}
\usepackage[shortlabels]{enumitem}
\newtheorem{lemma}{Lemma}
\newtheorem{theorem}{Theorem}
\newcommand\E{{\mathbb E}}
\newcommand\Pbb{{\mathbb P}}
\newcommand\Pcal{{\mathcal P}}
\newcommand\R{{\mathbb R}}
\newcommand\Z{{\mathbb Z}}
\newcommand\cA{{\mathsf C}}
\newcommand\ellpart{{\ell}}
\newcommand\one{{\mathbf 1}}
\title{Elliptic $b$-Hurwitz theory and Jack heat trace}
\author{Thibaut Lemoine}
\date{Source: arXiv:2609.12256v1 (CC BY 4.0)}
\begin{document}
\maketitle

\noindent\emph{Source and licence.} Elliptic $b$-Hurwitz theory and Jack heat trace. Version arXiv:2609.12256v1 (CC BY 4.0), distributed by arXiv under the Creative Commons Attribution 4.0 International licence, http://creativecommons.org/licenses/by/4.0/, as recorded in the arXiv licence metadata for this exact version and shown on the abstract page https://arxiv.org/abs/2609.12256v1. Full licence notice: \texttt{licenses/arxiv-2609.12256-CC-BY-4.0.txt}.\par\smallskip

\noindent\textit{Editorial scope.} Counted unit: the article's theorem thm:heat-expansion (source0.tex lines 490--496). The article's complete original proof of it is delivered with it (source0.tex lines 498--549, one \texttt{proof} environment of 52 lines). No mathematics is written, repaired or re-derived: the statement and the proof are the article's own lines, in the article's own notation, and no retained line is substituted or re-flowed.  Retained uncounted as the source-local context the proof needs: lines 308--310; lines 359--361; lines 431--433; lines 453--455; lines 458--460; lines 463--471.  Nothing was omitted from any retained span, so the package carries no omission record.  No retained line contains a citation, so the wrapper carries no bibliography and no bibliography entry.  No retained line contains a figure environment, an image-inclusion command or drawing code, so no image is delivered and the package declares no asset dependency.  Editorial formatting only, in addition: an AMS \texttt{amsart} preamble that restates the article's own declarations and macros used by the retained lines, namely the environments \texttt{lemma}, \texttt{theorem} on separate counters, together with the standard packages those lines need.  The retained environments print in this document as Lemma 1, Theorem 1 in order of appearance, whereas the article prints the same items with the numbering of its own full document.  The complete native source of the article as distributed in the arXiv e-print for this exact version is bundled unchanged as \texttt{sources/210092/source0.tex}.
\begin{equation}\label{eq:content-crude}
|\cA_\alpha(\lambda)|\leq \frac{\alpha+1}{2}|\lambda|(|\lambda|-1)\leq \frac{\alpha+1}{2}|\lambda|^2.
\end{equation}

\begin{equation}\label{eq:Ealpha}
E_\alpha(\mu,\nu,m)=|\mu|+|\nu|+\alpha m^2,\qquad F_\alpha(\mu,\nu,m)=\cA_\alpha(\mu)+\cA_\alpha(\nu)+\alpha m(|\mu|-|\nu|)-\frac{\alpha(\alpha-1)}2m^2.
\end{equation}

\begin{equation}\label{eq:discrete-Gaussian}
\Pbb_{\alpha,t}(M=m)=\frac{e^{-t\alpha m^2/2}}{\Theta_\alpha(t)}.
\end{equation}

\begin{equation}\label{eq:q-tail}
\sum_{\lambda\in\Pcal:|\lambda|>L}|\lambda|^s q^{|\lambda|} \leq Ce^{-cL},\qquad L\geq1,
\end{equation}

\begin{equation}\label{eq:gaussian-moments}
\sup_{(t,\alpha)\in K} \sum_{m\in\Z}(1+|m|)^s e^{-t\alpha m^2/4}<\infty,
\end{equation}

\begin{lemma}\label{lem:trace-tail}
Fix a compact set $K\subset(0,\infty)^2$ of $(t,\alpha)$-parameters, and let $0<\gamma<1$.  Then there exist $C,c>0$ such that, uniformly on $K$,
\begin{align}\label{eq:trace-tail}
&\sum_{\substack{\mu\in\Pcal_{\leq A_N},\nu\in\Pcal_{\leq B_N},m\in\Z\\
|\mu|+|\nu|>N^\gamma}}
e^{-\frac t2\kappa_{\alpha,N}(\lambda_N(\mu,\nu,m))}
\leq Ce^{-cN^\gamma}.
\end{align}
\end{lemma}

\begin{theorem}\label{thm:heat-expansion}
Let $p\in\Z_{\geq0}$.  For every $t>0$ and $\alpha>0$,
\begin{equation}\label{eq:heat-expansion}
Z_{\alpha,N}(t) =\frac{\Theta_\alpha(t)}{\phi(q_t)^2} \sum_{r=0}^{p} \frac{(-t)^r}{r!\bigl(N+\alpha-1\bigr)^r} \E\left[F_\alpha(\mu,\nu,M)^r\right] +O_{t,\alpha,p}(\bigl(N+\alpha-1\bigr)^{-p-1}).
\end{equation}
The estimate is locally uniform in $(t,\alpha)\in(0,\infty)^2$.
\end{theorem}

\begin{proof}
Fix a compact parameter set $K=[t_0,t_1]\times[\alpha_0,\alpha_1]\subset(0,\infty)^2$.  All constants below are uniform on $K$.  Set $D=N+\alpha-1$ and choose $\gamma=1/4$.  Let
\[
T_N=\{|\mu|+|\nu|\leq N^\gamma\}.
\]
For all sufficiently large $N$, $T_N\subset\Omega_N$, because
\[
\ellpart(\mu)\leq|\mu|\leq N^\gamma<A_N, \qquad \ellpart(\nu)\leq|\nu|\leq N^\gamma<B_N.
\]
By Lemma~\ref{lem:trace-tail}, the contribution of $\Omega_N\setminus T_N$ to the exact trace formula is $O(e^{-cN^\gamma})$ and therefore $O(N^{-m})$ for every $m$.

On $T_N$, the crude estimate \eqref{eq:content-crude} and \eqref{eq:Ealpha} give
\begin{equation}\label{eq:F-crude}
|F_\alpha(\mu,\nu,m)| \leq C_{\alpha_1}\left(S^2+|m|S+m^2+1\right), \qquad S=|\mu|+|\nu|.
\end{equation}
Since $S\leq N^{1/4}$ and $D\asymp N$ uniformly on $K$, the terms $S^2/D$ and $1/D$ are bounded.  Moreover, for $N$ large the coefficient of $m^2/D$ in \eqref{eq:F-crude} is at most $\alpha/8$, uniformly on $K$, while Young's inequality gives
\[
\frac{|m|S}{D}\leq \varepsilon m^2+C_\varepsilon\frac{S^2}{D^2}
\]
with $\varepsilon>0$ chosen uniformly small.  Hence, for all sufficiently large $N$,
\begin{equation}\label{eq:F-exp-control}
\frac{t}{D}|F_\alpha(\mu,\nu,m)| \leq \frac{t\alpha}{4}m^2+C_K.
\end{equation}

For $x\in\R$, Taylor's formula gives
\begin{equation}\label{eq:Taylor-bound}
\left|e^{-x}-\sum_{r=0}^{p}\frac{(-x)^r}{r!}\right| \leq \frac{|x|^{p+1}}{(p+1)!}e^{|x|}.
\end{equation}
Apply this with $x=tF_\alpha/D$ inside the exact probabilistic representation.  By \eqref{eq:F-exp-control}, the Gaussian factor from \eqref{eq:discrete-Gaussian} times $e^{|x|}$ is bounded by
\[
C_K e^{-t\alpha m^2/4}.
\]
The factor $|F_\alpha|^{p+1}$ is bounded by a fixed polynomial in $S$ and $|m|$ by \eqref{eq:F-crude}.  The summability bounds \eqref{eq:q-tail} and \eqref{eq:gaussian-moments} therefore imply
\[
\E\left[ |F_\alpha|^{p+1}e^{t|F_\alpha|/D}\one_{T_N} \right]\leq C_{K,p}.
\]
The Taylor remainder is thus $O(D^{-p-1})$.

We have proved
\begin{align*}
Z_{\alpha,N}(t)
&=\frac{\Theta_\alpha(t)}{\phi(q_t)^2}
\sum_{r=0}^{p}\frac{(-t)^r}{r!D^r}
\E\left[F_\alpha^r\one_{T_N}\right]
+O(D^{-p-1}).
\end{align*}
Finally, by \eqref{eq:F-crude} and the exponential tails of the product $q_t$-uniform law together with the discrete Gaussian law,
\[
\E\left[|F_\alpha|^r\one_{T_N^c}\right] =O(e^{-cN^\gamma})
\]
for each fixed $r$.  We may therefore remove $\one_{T_N}$ at an error smaller than any power of $N$.  This proves \eqref{eq:heat-expansion}, uniformly on $K$.
\end{proof}

\end{document}
